% Bibliografie

\newlength{\cslhangindent}
\setlength{\cslhangindent}{1.5em}

\newenvironment{CSLReferences}[2]
{%
	\begin{list}{}{%
			\setlength{\leftmargin}{\cslhangindent}%
			\setlength{\itemindent}{-\cslhangindent}%
			\setlength{\itemsep}{\baselineskip}%
			\setlength{\parsep}{0pt}%
		}%
		}
		{\end{list}}

\newcommand{\CSLLeftMargin}[1]{}
\newcommand{\CSLRightInline}[1]{#1}
\newcommand{\CSLIndent}[1]{\hspace{1.5em}}

% citeproc compatibility
\providecommand{\citeproc}[2]{#2}
\providecommand{\citeproctext}{}
\providecommand{\citeprocnum}[1]{#1}
\providecommand{\citeprocdate}[1]{#1}
\providecommand{\citeprocauthor}[1]{#1}

% Literaturverzeichnis: Klammern vor Labels unterdruecken
% (verhindert leere [] im thebibliography-Output)
\makeatletter
\renewcommand\@biblabel[1]{}
\makeatother
